\section{Parcial 2 -- 2015}

% =====================================================================
\subsection{Ejercicio 1 — Compartimentos de confidencialidad (restaurante)}
% =====================================================================

\begin{consigna}
Un famoso restaurante de 5 estrellas en la guía Michelin se destaca por todo lo que sirve: platos salados, platos dulces y bebidas (vinos y tragos). Se tienen los siguientes sujetos:
\begin{itemize}
  \item Jorge, José y Juan son los dueños.
  \item Manuel y Ramón son chefs.
  \item Luis es uno de los mozos.
  \item Andrea es una cliente.
\end{itemize}
Considerar los siguientes objetos:
\begin{itemize}
  \item Recetas (tanto de platos dulces, como de salados, como de tragos).
  \item Menú (se considera menú al listado de platos con sus respectivos precios).
  \item Listas de Stock (de productos para platos dulces, salados y tragos).
  \item Comanda (pedido que se hace al camarero-mozo).
\end{itemize}
Diseñar un modelo de compartimentos de confidencialidad, que permita garantizar los siguientes requisitos, justificando el cumplimiento de cada uno.
\begin{enumerate}
  \item Jorge puede leer todas las recetas, las listas de stock, el menú y las comandas.
  \item José, el somellier, es el único que escribe recetas de tragos. También lee stock de productos para tragos.
  \item Juan es el único que escribe recetas de platos (tanto dulces como salados). También lee stock de productos para platos.
  \item Los chefs pueden leer las recetas, pero no modificarlas.
  \item Los chefs pueden escribir en las listas de stocks para recomendar la compra de productos de su área.
  \item Ramón es el jefe de la barra, y sólo se ocupa de armar tragos. No cocina platos.
  \item Manuel dirige la cocina de platos dulces y salados.
  \item Los mozos sólo leen el menú y escriben comandas (no leen ni escriben nada más).
  \item Los clientes sólo leen menús.
\end{enumerate}
\end{consigna}

\paragraph{Temas.}
Clase 6 — Políticas y Control de Acceso (\S Bell--LaPadula; \S Compartimentos, reticulado y dominancia): modelo de confidencialidad por compartimentos, regla \emph{read-down / write-up}, relación de dominancia y \emph{need-to-know}.

\paragraph{Qué necesitás saber.}
\begin{itemize}
  \item Un \textbf{compartimento} es un par $(L, C)$: un nivel jerárquico $L$ y un conjunto de \textbf{categorías} (etiquetas) $C$. Aquí los niveles jerárquicos casi no importan; lo que ordena el ejercicio son las \textbf{categorías} (tragos / platos / menú / comanda).
  \item \textbf{Relación de dominancia:} $(L,C)\ \mathbf{dom}\ (L',C') \iff L'\le L \ \wedge\ C'\subseteq C$.
  \item \textbf{Reglas de Bell--LaPadula} (es un modelo de \emph{confidencialidad}, exactamente lo que pide la consigna):
  \begin{itemize}
    \item Seguridad simple — \emph{read down}: $s$ lee $o$ si $L(s)\ \mathbf{dom}\ L(o)$.
    \item Propiedad $\ast$ — \emph{write up}: $s$ escribe $o$ si $L(o)\ \mathbf{dom}\ L(s)$.
  \end{itemize}
  \item El reticulado deja pares \textbf{incomparables} (ninguna lista es subconjunto de la otra): quedan \textbf{efectivamente separados} — ni lectura ni escritura entre ellos. Es justo lo que se necesita para que la barra (tragos) y la cocina (platos) no se crucen.
\end{itemize}

\paragraph{Notas y conexiones.}
El restaurante es el clásico ejemplo de \emph{need-to-know}: ``un general a cargo de criptografía no tiene por qué conocer los secretos de los misiles''. Acá Ramón (barra) no necesita las recetas de platos y Manuel (cocina) no necesita las recetas de tragos. La compartimentalización tipo BLP, según la Clase 6, ``rara vez se aplica a todo un sistema'' y se reserva para separar dominios (finanzas/producción/marketing) — aquí, tragos/platos. \emph{Error clásico:} invertir read-down/write-up. Recordar: para que José sea el \textbf{único que escribe} recetas de tragos, esas recetas deben dominar (estar ``por encima de'') quienes solo las leen, y José debe estar en el mismo compartimento para poder escribir (write-up exige $L(o)\ \mathbf{dom}\ L(s)$; escribir-y-leer el mismo objeto exige misma clasificación, propiedad $\ast$-fuerte).

\paragraph{Resolución.}
Usamos un único nivel jerárquico base y modelamos todo con \textbf{categorías}. Definimos las categorías:
\[
  \mathsf{T}=\text{tragos},\quad \mathsf{P}=\text{platos},\quad \mathsf{M}=\text{menú},\quad \mathsf{K}=\text{comandas}.
\]
Las recetas/stock de tragos llevan $\{\mathsf{T}\}$; las de platos, $\{\mathsf{P}\}$; el menú, $\{\mathsf{M}\}$; las comandas, $\{\mathsf{K}\}$. Para forzar ``único que escribe'' separamos, dentro de cada rama, el objeto a escribir del objeto que otros solo leen, usando dos niveles ($\mathrm{alto}>\mathrm{bajo}$).

\medskip
\textbf{Etiquetas de objetos.}
\begin{center}
\renewcommand{\arraystretch}{1.2}
\begin{tabular}{ll}
\toprule
Objeto & Compartimento $(L,C)$ \\
\midrule
Recetas de tragos        & $(\mathrm{alto},\{\mathsf{T}\})$ \\
Stock de tragos          & $(\mathrm{bajo},\{\mathsf{T}\})$ \\
Recetas de platos        & $(\mathrm{alto},\{\mathsf{P}\})$ \\
Stock de platos          & $(\mathrm{bajo},\{\mathsf{P}\})$ \\
Menú                     & $(\mathrm{bajo},\{\mathsf{M}\})$ \\
Comanda                  & $(\mathrm{bajo},\{\mathsf{K}\})$ \\
\bottomrule
\end{tabular}
\end{center}

\textbf{Etiquetas de sujetos.}
\begin{center}
\renewcommand{\arraystretch}{1.2}
\begin{tabular}{lll}
\toprule
Sujeto & Compartimento $(L,C)$ & Rol \\
\midrule
Jorge (dueño)      & $(\mathrm{alto},\{\mathsf{T},\mathsf{P},\mathsf{M},\mathsf{K}\})$ & lee todo \\
José (somellier)   & $(\mathrm{alto},\{\mathsf{T}\})$ & escribe recetas tragos, lee stock tragos \\
Juan (dueño)       & $(\mathrm{alto},\{\mathsf{P}\})$ & escribe recetas platos, lee stock platos \\
Ramón (chef barra) & $(\mathrm{bajo},\{\mathsf{T}\})$ & lee recetas tragos, escribe stock tragos \\
Manuel (chef cocina)& $(\mathrm{bajo},\{\mathsf{P}\})$ & lee recetas platos, escribe stock platos \\
Luis (mozo)        & $(\mathrm{bajo},\{\mathsf{M},\mathsf{K}\})$ & lee menú, escribe comanda \\
Andrea (cliente)   & $(\mathrm{bajo},\{\mathsf{M}\})$ & lee menú \\
\bottomrule
\end{tabular}
\end{center}

\textbf{Verificación requisito por requisito} (read: $L(s)\,\mathbf{dom}\,L(o)$; write: $L(o)\,\mathbf{dom}\,L(s)$):
\begin{enumerate}
  \item \emph{Jorge lee todo.} $L(\text{Jorge})=(\mathrm{alto},\{\mathsf{T},\mathsf{P},\mathsf{M},\mathsf{K}\})$ domina a todos los objetos (nivel $\ge$ y todas las categorías $\subseteq$). $\Rightarrow$ lee recetas, stock, menú y comandas. \checkmark
  \item \emph{José único que escribe recetas de tragos.} write exige $(\mathrm{alto},\{\mathsf{T}\})\,\mathbf{dom}\,L(\text{José})=(\mathrm{alto},\{\mathsf{T}\})$: se cumple (misma etiqueta). Ramón está en $(\mathrm{bajo},\{\mathsf{T}\})$: para escribir las recetas necesitaría $(\mathrm{alto},\{\mathsf{T}\})\,\mathbf{dom}\,(\mathrm{bajo},\{\mathsf{T}\})$, que es \textbf{falso} ($\mathrm{alto}\not\le\mathrm{bajo}$) $\Rightarrow$ Ramón \textbf{no} escribe recetas. Lectura de stock de tragos por José: $(\mathrm{alto},\{\mathsf{T}\})\,\mathbf{dom}\,(\mathrm{bajo},\{\mathsf{T}\})$ \checkmark. \checkmark
  \item \emph{Juan único que escribe recetas de platos + lee stock platos.} Idéntico a José pero con $\{\mathsf{P}\}$. Manuel en $(\mathrm{bajo},\{\mathsf{P}\})$ no puede escribirlas (mismo argumento de nivel). \checkmark
  \item \emph{Chefs leen recetas pero no las modifican.} Ramón lee recetas de tragos: $(\mathrm{bajo},\{\mathsf{T}\})\,\mathbf{dom}\,(\mathrm{alto},\{\mathsf{T}\})$ es \textbf{falso} por nivel $\Rightarrow$ ¡no podría leer hacia arriba! Corrección: las recetas que los chefs leen se etiquetan en $(\mathrm{bajo},\{\mathsf{T}\})$ y la versión ``editable única'' del somellier/dueño es la de nivel alto que se \emph{publica} hacia abajo; en BLP estricto la lectura es \emph{read-down}, así que las recetas deben estar a nivel $\le$ del chef. Reasignamos: las recetas de tragos son $(\mathrm{bajo},\{\mathsf{T}\})$ y solo José (alto) escribe ``hacia arriba'' sobre ellas vía propiedad $\ast$ ($L(o)\,\mathbf{dom}\,L(s)$ requiere $o$ alto). Con recetas en $(\mathrm{bajo},\{\mathsf{T}\})$: Ramón \textbf{lee} ($(\mathrm{bajo},\{\mathsf{T}\})\,\mathbf{dom}\,(\mathrm{bajo},\{\mathsf{T}\})$ \checkmark) y \textbf{no escribe} (write exige $(\mathrm{bajo},\{\mathsf{T}\})\,\mathbf{dom}\,(\mathrm{bajo},\{\mathsf{T}\})$, que \emph{sí} se cumpliría). Para impedir la escritura del chef sin romper su lectura conviene poner las recetas \emph{por debajo} del chef: receta $=(\mathrm{bajo},\{\mathsf{T}\})$, chef $=(\mathrm{medio},\{\mathsf{T}\})$, somellier $=(\mathrm{bajo},\{\mathsf{T}\})$ — así el chef lee-down y solo el somellier (mismo nivel que el objeto) escribe. Ver tabla corregida abajo. \checkmark
  \item \emph{Chefs escriben en stock de su área.} write exige $L(\text{stock})\,\mathbf{dom}\,L(\text{chef})$. Con stock en el nivel más alto de la rama y chef debajo, write-up se cumple. \checkmark
  \item \emph{Ramón solo barra (tragos), no platos.} Ramón solo tiene categoría $\{\mathsf{T}\}$; los objetos de platos llevan $\{\mathsf{P}\}$ y $\{\mathsf{P}\}\not\subseteq\{\mathsf{T}\}$ $\Rightarrow$ \textbf{incomparable}: ni lee ni escribe nada de platos. \checkmark
  \item \emph{Manuel dirige cocina de platos.} Análogo a Ramón con $\{\mathsf{P}\}$; incomparable con $\{\mathsf{T}\}$. \checkmark
  \item \emph{Mozos: solo leen menú y escriben comandas.} Luis tiene $\{\mathsf{M},\mathsf{K}\}$. Lee menú $(\mathrm{bajo},\{\mathsf{M}\})$ \checkmark. Escribe comanda: write exige $(\mathrm{bajo},\{\mathsf{K}\})\,\mathbf{dom}\,L(\text{Luis})$; poniendo la comanda en el mismo nivel y categoría $\{\mathsf{K}\}\subseteq\{\mathsf{M},\mathsf{K}\}$, se cumple. No tiene $\mathsf{T}$ ni $\mathsf{P}$ $\Rightarrow$ nada de recetas/stock. \checkmark
  \item \emph{Clientes solo leen menú.} Andrea $=(\mathrm{bajo},\{\mathsf{M}\})$: lee menú \checkmark; sin $\mathsf{K}$ no escribe comandas; sin $\mathsf{T}/\mathsf{P}$ no toca recetas/stock. \checkmark
\end{enumerate}

\textbf{Asignación de niveles corregida} (para que ``chef lee receta pero no la escribe'' y ``único escritor'' convivan): dentro de cada rama de categoría usar tres niveles $\mathrm{stock}=\mathrm{alto} > \mathrm{chef}=\mathrm{medio} > \mathrm{receta}=\mathrm{bajo}$, con el escritor único (somellier/dueño) etiquetado al mismo nivel del objeto que escribe.
\begin{center}
\renewcommand{\arraystretch}{1.2}
\begin{tabular}{lll}
\toprule
Rama tragos & Compartimento & Justificación \\
\midrule
Receta tragos & $(1,\{\mathsf{T}\})$ & objeto que el chef lee-down \\
José (escribe receta) & $(1,\{\mathsf{T}\})$ & mismo nivel $\Rightarrow$ write \checkmark \\
Ramón (chef) & $(2,\{\mathsf{T}\})$ & lee receta (down) \checkmark; no la escribe ($1\not\ge 2$) \\
Stock tragos & $(3,\{\mathsf{T}\})$ & write-up del chef \checkmark; lee-down José/Jorge \\
\bottomrule
\end{tabular}
\end{center}
La rama platos es simétrica con $\{\mathsf{P}\}$ (Juan como escritor de recetas, Manuel como chef). Jorge se etiqueta en el nivel máximo con todas las categorías para leer-down todo. La separación tragos/platos queda garantizada por \textbf{incomparabilidad de categorías}, que es el corazón del reticulado BLP.

% =====================================================================
\subsection{Ejercicio 2 — Verdadero/Falso (pentesting, Anderson+salt, STRIDE)}
% =====================================================================

\begin{consigna}
Indica en cada caso si la afirmación es verdadera o falsa. Justificar en todos los casos. Si no se justifica correctamente, la respuesta quedará invalidada.
\begin{enumerate}
  \item En las pruebas de pentesting siempre se intenta explotar las vulnerabilidades encontradas.
  \item Si un atacante puede probar $G$ passwords por segundo, y las passwords están compuestas por símbolos en \texttt{[a-zA-Z0-9]}, entonces las passwords deberían tener 5 símbolos y 5 bits de SALT para que la probabilidad de adivinarlas en el lapso de un año sea menor a 0,5. \textbf{(El valor de $G$ aparecía como \texttt{EMBED Equation.3} en el original; no es recuperable.)}
  \item STRIDE forma parte del método de Modelado de Amenazas de Microsoft.
\end{enumerate}
\end{consigna}

\paragraph{Temas.}
Clase 10b — Pentesting (\S Metodología de Hipótesis de Falla, ``explotar es el último recurso''); Clase 7 — Autenticación (\S Fórmula de Anderson; \S Salting); Clase 8 — Principios de Diseño y Vulnerabilidades (\S Modelo STRIDE / Threat Modeling de Microsoft).

\paragraph{Qué necesitás saber.}
\begin{itemize}
  \item \textbf{Pentesting:} la fase de Prueba prioriza por nivel de acceso y \textbf{explotar es el último recurso} (mejor analizar documentación/comportamiento); la prueba debe ser repetible, conclusiva y \emph{lo menos intrusiva posible} (un exploit puede causar DoS).
  \item \textbf{Anderson:} $P = T\,G/N$, con $P$ probabilidad de adivinar, $T$ tiempo, $G$ pruebas/seg, $N$ tamaño del espacio. El \textbf{salt no entra en $N$}: no agranda el espacio de una clave concreta; lo que hace es impedir la \textbf{paralelización} del ataque (rainbow tables / un mismo hash para dos usuarios). Por eso ``$5$ bits de SALT'' no cambia $P$ de adivinar \emph{una} clave puntual.
  \item \textbf{STRIDE:} las 6 categorías de amenazas (Spoofing, Tampering, Repudiation, Information Disclosure, DoS, Elevation of Privilege) son parte del \textbf{Threat Modeling de Microsoft}, combinadas con un DFD.
\end{itemize}

\paragraph{Notas y conexiones.}
El ítem (a) es la \textbf{trampa recurrente} de pentesting: ``explotar es el último recurso''. El ítem (b) mezcla Anderson (Clase 7) con un distractor de salt; el alfabeto \texttt{[a-zA-Z0-9]} tiene $26+26+10=62$ símbolos. El ítem (c) es teoría directa de Clase 8.

\paragraph{Resolución.}
\textbf{(a) FALSO.} En pentesting \emph{no} siempre se explota: explotar es el \textbf{último recurso} porque puede ser intrusivo (incluso provocar un DoS). Se prefiere analizar documentación y comportamiento; la prueba debe ser repetible, conclusiva y mínimamente intrusiva.

\textbf{(b) FALSO} (con dos argumentos, uno independiente de $G$). Espacio de claves con $5$ símbolos sobre $62$:
\begin{equation*}
  N = 62^{5} = 916\,132\,832 \approx 9.16\times10^{8}.
\end{equation*}
Aplicando Anderson sobre un año $T = 365\cdot24\cdot3600 = 31\,536\,000$ s, la condición $P<0.5$ exige
\begin{equation*}
  P=\frac{T\,G}{N}<0.5 \;\Longleftrightarrow\; G < \frac{0.5\,N}{T} = \frac{0.5\cdot 62^{5}}{31\,536\,000} \approx 14.5\ \text{pruebas/s}.
\end{equation*}
O sea, $5$ símbolos solo bastan si el atacante prueba menos de $\sim 15$ claves por segundo — un valor irrisorio para cualquier ataque offline real (millones/seg). Salvo que el $G$ embebido fuera ridículamente bajo, la afirmación es \textbf{falsa}. Argumento \textbf{independiente de $G$}: los \textbf{5 bits de SALT no influyen} en la probabilidad de adivinar \emph{una} contraseña concreta —el salt no agranda $N$, solo impide la paralelización/precómputo (rainbow tables) y que dos claves iguales den el mismo hash—. Por lo tanto agregar ``5 bits de SALT'' como condición para bajar $P$ es conceptualmente incorrecto. \textbf{Falso}. \emph{(El $G$ original es \texttt{EMBED Equation.3} no recuperable; la conclusión no depende de su valor por el argumento del salt.)}

\textbf{(c) VERDADERO.} STRIDE es el modelo de 6 categorías de amenazas de \textbf{Microsoft}, parte de su metodología de Threat Modeling, que se cruza con un DFD para identificar al menos dos amenazas por categoría.

% =====================================================================
\subsection{Ejercicio 3 — Secreto compartido de Shamir (3,4) + entropía}
% =====================================================================

\begin{consigna}
Considerar el conjunto $C$ de sombras para un esquema de secreto compartido de Shamir $(k,n)$, donde $k=3$ y $n=4$ sobre $\mathbb{Z}_p$ \textbf{(el cuerpo aparecía como \texttt{EMBED Equation.3}; ver más abajo el cuerpo asumido)}.
\[
  C = \{(1,6),(2,1),(3,0),(4,3)\}.
\]
\begin{enumerate}
  \item Obtener el secreto.
  \item Obtener el polinomio utilizado para obtener el conjunto de sombras.
  \item Expresar el significado de un esquema de secreto umbral $(3,n)$ utilizando el concepto de entropía.
\end{enumerate}
\end{consigna}

\paragraph{Temas.}
Clase 8 — Principios de Diseño (\S Flujo de información: introducción a Shamir y entropía) y Clase 9 — Flujo de Información y Malware (\S Entropía de Shannon, $H(X)$ y $H(X\mid Y)$). Secreto compartido de Shamir: interpolación de Lagrange sobre un cuerpo finito.

\paragraph{Qué necesitás saber.}
\begin{itemize}
  \item \textbf{Shamir $(k,n)$:} el secreto $s=f(0)$ de un polinomio de grado $k-1$ sobre $\mathbb{Z}_p$. Con $k$ sombras $(x_i,y_i)$ se reconstruye $f$ por \textbf{interpolación de Lagrange}; con menos de $k$ no se obtiene \emph{ninguna} información.
  \item Acá $k=3$ $\Rightarrow$ polinomio de grado $2$: $f(x)=a_2x^2+a_1x+a_0$, secreto $a_0=f(0)$.
  \item \textbf{Cuerpo asumido:} la consigna perdió $p$. Las cuatro sombras tienen valores en $\{0,\dots,6\}$ y son \emph{consistentes} con un único polinomio de grado $2$ sobre $\mathbb{Z}_7$ (el menor primo que contiene todos los valores); por eso tomamos $\boxed{p=7}$.
  \item \textbf{Entropía de Shannon:} $H(X)=-\sum_i p(x_i)\lg p(x_i)$; condicional $H(X\mid Y)=\sum_j p(y_j)H(X\mid Y=y_j)$. Mide la \textbf{incertidumbre}; máxima $\lg n$ (uniforme), mínima $0$ (certeza).
\end{itemize}

\paragraph{Notas y conexiones.}
Tener $n=4$ sombras con $k=3$ hace el sistema \textbf{sobredeterminado}: alcanzan 3 para reconstruir y la 4ta sirve de verificación de consistencia (que es como deducimos el cuerpo). El inciso (c) conecta Shamir con la entropía vista en Clase 9: ``saber menos de $k$ sombras = entropía máxima sobre el secreto''.

\paragraph{Resolución.}
\textbf{(b) Polinomio.} Buscamos $f(x)=a_2x^2+a_1x+a_0$ sobre $\mathbb{Z}_7$ con $f(1)=6,\ f(2)=1,\ f(3)=0$:
\begin{align*}
  a_0+a_1+a_2 &\equiv 6 \pmod 7,\\
  a_0+2a_1+4a_2 &\equiv 1 \pmod 7,\\
  a_0+3a_1+2a_2 &\equiv 0 \pmod 7 \quad (\text{pues } 9\equiv 2).
\end{align*}
Restando la 1ra a la 2da: $a_1+3a_2\equiv 1-6\equiv 2 \pmod 7$. Restando la 1ra a la 3ra: $2a_1+a_2\equiv 0-6\equiv 1 \pmod 7$. De la primera, $a_1\equiv 2-3a_2$; sustituyendo: $2(2-3a_2)+a_2\equiv 1 \Rightarrow 4-5a_2\equiv 1 \Rightarrow -5a_2\equiv -3 \Rightarrow 2a_2\equiv 4 \Rightarrow a_2\equiv 2$. Entonces $a_1\equiv 2-6\equiv -4\equiv 3$ y $a_0\equiv 6-a_1-a_2 = 6-3-2 = 1$. Resulta
\begin{equation*}
  \boxed{f(x) = 2x^2 + 3x + 1 \pmod 7}.
\end{equation*}
\emph{Verificación (incluida la 4ta sombra):} $f(1)=2+3+1=6$; $f(2)=8+6+1=15\equiv 1$; $f(3)=18+9+1=28\equiv 0$; $f(4)=32+12+1=45\equiv 3$. Las cuatro sombras coinciden, lo que confirma $p=7$.

\textbf{(a) Secreto.} El secreto es el término independiente:
\begin{equation*}
  s = f(0) = a_0 = \boxed{1}.
\end{equation*}
Equivalentemente, por Lagrange sobre las tres primeras sombras: $f(0)=\sum_i y_i\prod_{j\neq i}\frac{-x_j}{x_i-x_j}\pmod 7 = 1$.

\textbf{(c) Significado con entropía.} Sea $S$ la variable aleatoria del secreto. Un esquema umbral $(3,n)$ significa, en términos de entropía:
\begin{itemize}
  \item Con \textbf{menos de 3} sombras no se obtiene \emph{ninguna} información sobre el secreto: la incertidumbre permanece \textbf{máxima}, $H(S\mid \text{cualquier conjunto de }<3\text{ sombras}) = H(S) = \lg|\mathbb{Z}_p|$ (uniforme). Conocer $0,1$ o $2$ sombras no reduce la entropía.
  \item Con \textbf{3 o más} sombras la incertidumbre cae a \textbf{cero}: $H(S\mid \ge 3\text{ sombras}) = 0$ (el secreto queda determinado).
\end{itemize}
Es decir, hay \textbf{flujo de información} sobre $S$ recién al alcanzar el umbral de $k=3$ sombras: por debajo del umbral $H$ no baja (no hay flujo), en el umbral $H$ salta de $\lg p$ a $0$.

% =====================================================================
\subsection{Ejercicio 4 — Resolución de conflictos en ACL + DMZ}
% =====================================================================

\begin{consigna}
En cada caso selecciona la opción correcta. Justificar en todos los casos. Si no se justifica correctamente, la respuesta quedará invalidada.

\medskip
\textbf{(a)} Si se tiene una lista de control de acceso $\mathrm{ACL}(o) = \{(\text{Manuel}, r), (\text{Directivos}, w), (\ast, w)\}$. Asumiendo que Manuel es Directivo:
\begin{enumerate}
  \item Manuel puede leer y escribir en $o$ si la política es \emph{first rule}.
  \item Manuel puede leer y escribir en $o$ si la política es \emph{grant all}.
  \item Manuel puede leer y escribir en $o$ si la política es \emph{override}.
  \item Manuel puede leer y escribir en $o$ si la política es \emph{augment}.
\end{enumerate}

\textbf{(b)} En una red bien diseñada, la zona demilitarizada (DMZ) nunca debería contener:
\begin{enumerate}
  \item Servidores de email.
  \item Servidores de DNS.
  \item Servidores de subred de desarrollo.
  \item Servidores de Web Proxy.
\end{enumerate}
\end{consigna}

\paragraph{Temas.}
Clase 6 — Políticas y Control de Acceso (\S Conflictos en ACL; \S Derechos por defecto: override/augment); Clase 10 — Seguridad en la Empresa (\S Defense-in-depth; DMZ).

\paragraph{Qué necesitás saber.}
\begin{itemize}
  \item \textbf{Conflictos de ACL} (Clase 6) cuando varias entradas aplican a un sujeto:
  \begin{itemize}
    \item \emph{Permitir (unión):} permite si \emph{alguna} entrada lo da.
    \item \emph{Grant All (intersección):} deniega si \emph{alguna} entrada lo deniega — exige que \textbf{todas} den el acceso.
    \item \emph{First Rule / First Match:} aplica la \textbf{primera} entrada que matchea.
    \item \emph{Override / Best match:} si hay entrada específica del sujeto, se usa esa; el default $\ast$ solo si no la hay.
    \item \emph{Augment:} se parte del default $\ast$ y se le \textbf{suman} las entradas explícitas.
  \end{itemize}
  \item \textbf{DMZ} (Clase 10): zona expuesta entre el firewall externo y el interno; aloja servicios que \emph{deben} ser accesibles desde Internet (web, correo, DNS público, proxy). Lo interno/sensible va detrás del firewall interno (LAN).
\end{itemize}

\paragraph{Notas y conexiones.}
El inciso (a) es el ejercicio resuelto de ACL de la Clase 6 trasladado a este enunciado: hay que evaluar las cuatro políticas y ver con cuál Manuel obtiene \textbf{ambos} permisos ($r$ \emph{y} $w$). Las entradas que aplican a Manuel: $(\text{Manuel},r)$ específica, $(\text{Directivos},w)$ de grupo, $(\ast,w)$ default.

\paragraph{Resolución.}
\textbf{(a)} Entradas que aplican a Manuel: específica $r$, grupo Directivos $w$, default $\ast$ $w$.
\begin{center}
\renewcommand{\arraystretch}{1.3}
\begin{tabular}{>{\bfseries}l l l}
\toprule
Política & Resultado para Manuel & ¿$r$ y $w$? \\
\midrule
First rule & primera entrada $(\text{Manuel},r)\Rightarrow$ solo $r$ & No \\
Grant all  & intersección $r\cap w\cap w = \varnothing$ & No \\
Override   & usa la específica $(\text{Manuel},r)\Rightarrow$ solo $r$ & No \\
Augment    & default/grupo $w$ $+$ específica $r$ $\Rightarrow r,w$ & \textbf{Sí} \\
\bottomrule
\end{tabular}
\end{center}
\textbf{Correcta: opción 4 (\emph{augment}).} Con \emph{augment} se parte de lo que dan el grupo Directivos y el default ($w$) y se le agrega la entrada específica de Manuel ($r$): Manuel obtiene $\{r,w\}$. Con \emph{first rule} y \emph{override} domina su entrada específica $(r)$ y pierde el $w$; con \emph{grant all} la intersección de $r,w,w$ es vacía (ni siquiera $w$, porque la entrada específica no incluye $w$).

\textbf{(b)} \textbf{Correcta: opción 3 (servidores de subred de desarrollo).} La DMZ aloja exactamente los servicios que deben ser alcanzables desde Internet: \emph{web}, \emph{correo}, \emph{DNS público} y \emph{web proxy} (las opciones 1, 2 y 4 \textbf{sí} pertenecen a la DMZ). Una \textbf{subred de desarrollo} es un activo interno, no expuesto: debe vivir detrás del firewall interno, en la LAN. Colocar desarrollo en la DMZ expondría sistemas sensibles y no endurecidos a la red externa, violando la lógica de la zona desmilitarizada.
